\section{Exact histories, a fixed clique palette, and hardness}
\label{sec:circuit}

Let
\[
G_2=\{X,\CX,\CCX,H\otimes H\}.
\]
Following Rudolph's gate-dependent definition
\cite[Definition~2.2]{rudolph2026universal}, \(\BQPone^{G_2}\) consists of
promise problems decided by a polynomial-time uniform family of exact
\(G_2\) circuits on an all-zero input, with one computational-basis output
bit, perfect completeness, and inverse-polynomial rejection probability in
the no case.  Rudolph's exact error reduction
\cite[Theorem~2.3]{rudolph2026universal} preserves perfect completeness and
stays in any gate set containing \(G_2\); we apply it first so that soundness
is at most \(1/3\).  The superscript is essential: no approximate compilation
or gate-independent identification is used below.

\begin{theorem}[Fixed-eight hardness]
\label{thm:weighted-hardness}
There is a deterministic polynomial-time many-one reduction from
\(\BQPone^{G_2}\) to the decision version of weighted gapped true normalized
persistence.  Every output satisfies
\[
\beta_d(\Xin)=8,\qquad
\beta_d(\Xin\to\Xout)=
\begin{cases}6,&x\in L,\\1,&x\notin L,\end{cases}
\]
and both endpoint degree-\(d\) Laplacians have inverse-polynomial gaps above
their complete kernels.  Hence the normalized values are \(3/4\) and \(1/8\).
Additive error \(1/24\), with decision threshold \(7/16\), distinguishes the
two cases with success probability at least \(2/3\).
\end{theorem}

The remainder of this section proves the theorem by combining a rational
history construction, a certified finite clique palette, and an exact
three-bit source combiner.

\subsection{Rational histories and the certified palette}

Write \(K=\sqrt2H\).  Replace each \(H\otimes H\) transition by a gain edge
\(K\) on one work bit followed by a loss edge \(K^{-1}=H/\sqrt2\) on the
other.  Their composite is exactly \(H\otimes H\), and each edge is enforced
by two orthonormal rational three-term vectors.  The vectors and their direct
calculation appear in \cref{app:rational-split}.  This construction is due to
Rudolph~\cite{rudolph2026universal}.

The propagation constraints for \(X,\CX,\CCX\) decompose, on their
constant-size active support, into two-time vectors
\[
(\ket{0,z}-\ket{1,Uz})/\sqrt2.
\]
Clock, input, and output penalties are basis projectors.  Unaffected work
qubits are identities and are implemented by the harmonic padding of
\cref{lem:scaling-padding}; they are not expanded into exponentially many
rank-one terms.

\begin{lemma}[Certified fixed palette]
\label{lem:palette}
Every local constraint in the split \(G_2\) history is implemented by a
member of one fixed family of weighted clique gadgets of total support
locality at most six.  Every member satisfies
\cref{ass:local-certificate,eq:interface}, with uniform constants, and the
graph construction is polynomial in the circuit size.
\end{lemma}

\begin{proof}
See \cref{app:palette}.  It records the finite integer certificates, exact
relabelings, selected-cycle guard closure, and polynomial-size bookkeeping.
\end{proof}

\subsection{Complete history kernels}

Let \(L\) be the number of legal clock positions after splitting the
Hadamard transitions.  Put \(s_t=\sqrt2\) immediately after a gain edge and
\(s_t=1\) otherwise.  If \(V_t\) is the unitary prefix after removing these
scalars, then
\begin{equation}
\label{eq:history-isometry}
\cV v=Z^{-1/2}\sum_{t=0}^{L-1}s_t\ket tV_tv,\qquad
Z=\sum_ts_t^2,\qquad L\le Z\le2L.
\end{equation}
Every gain edge is immediately followed by its matching loss edge, with no
gate between them, so \(s_0=s_{L-1}=1\).

\begin{proposition}[History kernels and gaps]
\label{prop:history}
Let \(C\) be the valid clean-input space, \(D=\dim C>0\), and let
\(M=J^*U^*P_{\rm acc}UJ\) be the compressed acceptance operator.  Put
\(S=\ker(I-M)\).  The initial history Hamiltonian has complete kernel
\(\cV C\), while adding the final rejection projector gives complete kernel
\(\cV S\).  If \(M|_{S^\perp}\preceq I/3\), then
\begin{equation}
\label{eq:history-gaps}
g_{\rm in}\ge\frac1{8L^2},\qquad
g_{\rm out}\ge\frac1{3L(8L^2+1)}\ge\frac1{27L^3}.
\end{equation}
The second statement includes \(S=0\).
\end{proposition}

The proof in \cref{app:source-details} uses a weighted path Poincar\'e
inequality.  The output compression is exactly
\begin{equation}
\label{eq:output-compression}
\cV^*H_{\rm rej}\cV=(I-M)/Z,
\end{equation}
and a two-positive-operator estimate turns this angle penalty into the second
gap in \cref{eq:history-gaps}.  Combining
\cref{thm:transfer,thm:quotient,lem:palette,prop:history} yields weighted
clique complexes \(\Xin\subseteq\Xout\) with
\begin{equation}
\label{eq:generic-circuit-transfer}
\beta_d(\Xin)=D,\qquad
\beta_d(\Xout)=\beta_d(\Xin\to\Xout)=\dim S,
\end{equation}
and inverse-polynomial gaps above both complete kernels.

\subsection{A fixed-eight source and proof of \cref{thm:weighted-hardness}}

Given a \(\BQPone^{G_2}\) verifier \(Q_x\) with acceptance probability \(p_x\),
add a maximally mixed label \(z\in\{0,1\}^3\), run \(Q_x\) unconditionally,
and accept according to
\[
[z=000]\ \lor\
\bigl(q\land[z\in\{001,010,011,100,101\}]\bigr),
\]
where \(q\) is the output of \(Q_x\).  The reversible predicate uses only
\(X,\CX,\CCX\), preserves \(z\), and uncomputes its scratch.  Its compressed
acceptance operator on the label is exactly
\begin{equation}
\label{eq:eight-label}
M_x=\operatorname{diag}(1,p_x,p_x,p_x,p_x,p_x,0,0).
\end{equation}
Thus the perfect eigenspace has dimension six when \(p_x=1\) and dimension
one when \(p_x\le1/3\); every off-perfect eigenvalue in the latter case is at
most \(1/3\).

\begin{proof}[Proof of \cref{thm:weighted-hardness}]
Apply \cref{eq:generic-circuit-transfer} to \cref{eq:eight-label}.  The two
ratios differ by \(5/8\), and \(1/24<5/16\).  The logical gaps in
\cref{eq:history-gaps} are explicit rationals.  The finite certificate data
supply explicit positive rational constants \(\bar c_1,\bar c_2\) small enough
that \cref{eq:lambda-E-choice} remains valid with \(c_1,c_2\) replaced by
\(\bar c_1,\bar c_2\).  Fix \(\eta_0=1/10\).  From the polynomially bounded
\(t_{\max}\), choose the least integer \(b\ge0\) for which
\(\lambda=2^{-b}\) satisfies the first inequality in
\cref{eq:lambda-E-choice} with \(\bar c_1\).

For endpoint \(i\), put
\[
r_i=\frac{\bar c_2}{2}\eta_0^2g_i\lambda^{26}
\]
and choose the least integer \(s_i\ge0\) such that \(2^{-s_i}\le r_i\).
Output the dyadic promise bound \(E_i=2^{-s_i}\).  Then
\(0<E_i\le r_i<c_2\eta_0^2g_i\lambda^{26}\).  Exact integer comparison of a
power of two with the rational \(r_i\) computes \(s_i\); since \(g_i^{-1}\)
and \(t_{\max}\) are polynomially bounded, \(b,s_i\) and the binary encodings
of \(\lambda,E_i\) have polynomial length.

Let \(n\) be the number of logical register factors in the split history.
The weighted graph has \(O(n+t_{\max})\) vertices and at most
\(O(n^2+t_{\max}n)\) explicitly generated edges, because the register and
every local interior are explicit and locality is constant.  Hence the map
outputs the two graphs, their common weight function, \(d,E_1,E_2\), and the
fixed decision threshold in deterministic polynomial time.
\end{proof}

For every fixed finite gate family covered by Rudolph's
Theorem~3.4~\cite{rudolph2026universal}, the exact simulation into \(G_2\) may
be composed with this reduction.  This is a separate gate-family-relative
statement, not gate-independent \(\BQPone\)-hardness.

The exact trace-versus-perfect-space obstruction and the mixed-spectator
extension to individual real Hadamards are recorded in
\cref{app:source-boundaries}.  Neither is needed for
\cref{thm:weighted-hardness}.
